% Part 3 | Status: todo | Original lines: 2449-2458
%!TEX root = ../../main.tex
\clearpage
% =====================================================================
\part{Integration and the theorems of Stokes and Gauss}\label{part:integration}
% =====================================================================

A covector eats a vector; a curve provides, at every instant, a vector (its velocity). Integrating a $1$-form along a curve is therefore
nothing but contracting it with the velocity and adding up. The same idea with a tangent $k$-frame integrates $k$-forms, and the
fundamental theorem of calculus then generalises to a single statement, Stokes' theorem, from which the divergence theorem follows by
a partial contraction.
% Section 16: Integration of $1$-forms | Part: 3 | Depends on: sec:fields, sec:tangent | Status: todo | Original lines: 2459-2462
%!TEX root = ../../../main.tex
% =====================================================================
\section{Integration of $1$-forms}\label{sec:int1}
% =====================================================================
% Section 16.1: Definition: integration is contraction with the velocity | Part: 3 | Depends on: sec:fields, sec:tangent | Status: todo | Original lines: 2463-2494
%!TEX root = ../../../main.tex
\subsection{Definition: integration is contraction with the velocity}\label{sec:int1def}

Let $\gamma:[a,b]\to M$ be a piecewise $C^1$ curve, with velocity $\dot\gamma(t)\in T_{\gamma(t)}M$.

\begin{definition}[Line integral]
The integral of a $1$-form $\omega$ along $\gamma$ is
\begin{equation}\label{eq:lineint}
\int_\gamma\omega:=\int_a^b\bigl\langle\omega_{\gamma(t)},\dot\gamma(t)\bigr\rangle\,\dd t=\int_a^b\omega_i(\gamma(t))\,\dot\gamma^i(t)\,\dd t .
\end{equation}
Equivalently, $\int_\gamma\omega=\int_{[a,b]}\gamma\pb\omega$, where the pull-back $\gamma\pb\omega=\omega(\dot\gamma)\,\dd t$ is a $1$-form on the interval.
\end{definition}

So the integrand is nothing but the contraction of the covector with the velocity vector, followed by an ordinary Riemann integral over the parameter.
% REORDERED: the general definition eq:pullback moved to Section sec:pushpull (merged with the GR definition there).
Recall from Section~\ref{sec:pushpull} that the pull-back $\phi\pb\omega$ of a form by a smooth map $\phi$ is defined by contraction with pushed-forward vectors, \eqref{eq:pullback}.
\emph{Why $\gamma\pb\omega=\omega(\dot\gamma)\,\dd t$ (steps).} Use the coordinates $\gamma^i(t)=x^i(\gamma(t))$. (1) The pull-back of the coordinate function $x^i$ is $x^i\circ\gamma=\gamma^i$, so $\gamma\pb\dd x^i=\dd(\gamma^i)=\dot\gamma^i\dd t$ (pull-back commutes with $\dd$).
(2) Pull-back is linear and multiplies the coefficient functions by composition: $\gamma\pb(\omega_i\dd x^i)=\omega_i(\gamma(t))\dot\gamma^i(t)\,\dd t$. (3) This is $\langle\omega_{\gamma(t)},\dot\gamma(t)\rangle\dd t$, as claimed.

\begin{example}[A line integral, step by step]\label{ex:lineint}
Compute $\int_\gamma\omega$ for $\omega=x\,\dd y$ and the parabola $\gamma(t)=(t,t^2)$, $0\le t\le1$.
\begin{enumerate}
\item Velocity: $\dot\gamma(t)=(1,2t)$.
\item Coefficients of $\omega$ on the curve: $\omega=\omega_x\dd x+\omega_y\dd y$ with $\omega_x=0$, $\omega_y=x$; on $\gamma$, $\omega_y=t$.
\item Contraction: $\langle\omega_{\gamma(t)},\dot\gamma(t)\rangle=0\cdot1+t\cdot2t=2t^2$.
\item Integral: $\int_0^12t^2\,\dd t=\frac23$.
\end{enumerate}
\emph{Check with another parametrisation of the same curve:} $\tilde\gamma(s)=(s^2,s^4)$, $0\le s\le1$. Then $\dot{\tilde\gamma}=(2s,4s^3)$, $\omega_y=s^2$, contraction $=s^2\cdot4s^3=4s^5$, and $\int_0^14s^5\,\dd s=\frac23$: the same.
\emph{Reversed curve:} $\hat\gamma(t)=(1-t,(1-t)^2)$ gives $\dot{\hat\gamma}=(-1,-2(1-t))$, contraction $=-2(1-t)^2$ and the integral is $-\frac23$: the sign flips (Proposition~\ref{prop:reparam}).
\end{example}

% Section 16.2: Why a covector? Riemann sums and infinitesimal displacements | Part: 3 | Depends on: sec:fields, sec:tangent | Status: todo | Original lines: 2495-2504
%!TEX root = ../../../main.tex
\subsection{Why a covector? Riemann sums and infinitesimal displacements}\label{sec:riemannsums}

Partition $[a,b]$ by $t_0<\cdots<t_N$ and put $\Delta\gamma_j=\gamma(t_{j+1})-\gamma(t_j)$ (in a chart). Then
\[
\int_\gamma\omega=\lim_{\max|\Delta t_j|\to0}\sum_j\bigl\langle\omega_{\gamma(t_j)},\Delta\gamma_j\bigr\rangle .
\]
Why does the sum converge to \eqref{eq:lineint}? On a short interval $\Delta\gamma_j=\dot\gamma(t_j)\Delta t_j+o(\Delta t_j)$ (definition of the derivative). By linearity of the contraction,
$\langle\omega_{\gamma(t_j)},\Delta\gamma_j\rangle=\langle\omega_{\gamma(t_j)},\dot\gamma(t_j)\rangle\Delta t_j+o(\Delta t_j)$, so the sum is, up to a vanishing error, a Riemann sum of the function $t\mapsto\langle\omega_{\gamma(t)},\dot\gamma(t)\rangle$.
The picture is Figure~\ref{fig:sheets}: each small displacement $\Delta\gamma_j$ pierces a certain number of ``sheets'' of $\omega_{\gamma(t_j)}$, and the line integral counts the total number of sheets crossed along the path. A $1$-form is therefore precisely the object which assigns a number to every infinitesimal oriented displacement, linearly. This also explains why $\int_\gamma\omega$ needs no metric: no length or angle is used, only the pairing \eqref{eq:pairing}.

% Section 16.3: The role of linearity | Part: 3 | Depends on: sec:fields, sec:tangent | Status: todo | Original lines: 2505-2522
%!TEX root = ../../../main.tex
\subsection{The role of linearity}\label{sec:linearity}

\begin{proposition}[Reparametrisation and orientation]\label{prop:reparam}
Let $\tau:[c,d]\to[a,b]$ be a $C^1$ diffeomorphism. If $\tau'>0$ then $\int_{\gamma\circ\tau}\omega=\int_\gamma\omega$; if $\tau'<0$ then $\int_{\gamma\circ\tau}\omega=-\int_\gamma\omega$.
\end{proposition}
\begin{proof}
\Step{1}{velocity of the new curve} By the chain rule, $\frac{\dd}{\dd s}(\gamma\circ\tau)(s)=\tau'(s)\,\dot\gamma(\tau(s))$.
\Step{2}{pull the scalar out} Because $\omega_p$ is linear, $\langle\omega,\tau'\dot\gamma\rangle=\tau'\langle\omega,\dot\gamma\rangle$. Hence $\int_{\gamma\circ\tau}\omega=\int_c^d\langle\omega_{\gamma(\tau(s))},\dot\gamma(\tau(s))\rangle\tau'(s)\,\dd s$.
\Step{3}{substitute} $t=\tau(s)$, $\dd t=\tau'(s)\,\dd s$. The integral becomes $\int_{\tau(c)}^{\tau(d)}\langle\omega_{\gamma(t)},\dot\gamma(t)\rangle\,\dd t$.
\Step{4}{limits} If $\tau'>0$, then $\tau(c)=a$ and $\tau(d)=b$, so we obtain $\int_\gamma\omega$. If $\tau'<0$, then $\tau(c)=b$ and $\tau(d)=a$; exchanging the limits gives $-\int_\gamma\omega$.
(Example~\ref{ex:lineint} is a numerical instance with $\tau(s)=s^2$ and with $\tau(t)=1-t$.)
\end{proof}

\begin{keyidea}
Linearity of the contraction in the velocity is exactly what makes line integrals independent of the parametrisation. Replacing $\omega(\dot\gamma)$ by a non-linear function of the velocity -- \eg the speed $|\dot\gamma|$ -- gives arc length,
which is also reparametrisation-invariant but not sign-reversing under change of orientation, and needs a metric. $1$-forms are the integrands of \emph{oriented} curves.
\end{keyidea}

% Section 16.4: Naturality | Part: 3 | Depends on: sec:fields, sec:tangent | Status: done | Original lines: 2523-2537
%!TEX root = ../../../main.tex
\subsection{Naturality}\label{sec:naturality}

For a smooth map $\phi:M\to N$, a curve $\gamma$ in $M$ and a $1$-form $\omega$ on $N$:
\begin{equation}\label{eq:natural}
\int_{\phi\circ\gamma}\omega=\int_\gamma\phi\pb\omega .
\end{equation}
\emph{Proof in steps:}
\begin{align*}
\int_{\phi\circ\gamma}\omega&=\int_a^b\Bigl\langle\omega_{\phi(\gamma(t))},\frac{\dd}{\dd t}(\phi\circ\gamma)(t)\Bigr\rangle\dd t &&\text{(definition \eqref{eq:lineint})}\\
&=\int_a^b\bigl\langle\omega_{\phi(\gamma(t))},\dd\phi_{\gamma(t)}\dot\gamma(t)\bigr\rangle\dd t &&\text{(chain rule)}\\
&=\int_a^b\bigl\langle(\phi\pb\omega)_{\gamma(t)},\dot\gamma(t)\bigr\rangle\dd t &&\text{(definition \eqref{eq:pullback})}\\
&=\int_\gamma\phi\pb\omega .
\end{align*}
In particular, line integrals of $1$-forms are invariant under changes of coordinates (take $\phi$ a diffeomorphism). In contrast, the classical ``work integral'' $\int\mathbf F\cdot\dd\mathbf r$ is invariant only because it is secretly $\int F^\flat$, where $F^\flat=g(F,\cdot)$ is the covector obtained from $F$ with the Euclidean metric: $F\cdot\dot r=g(F,\dot r)=\langle F^\flat,\dot r\rangle$.

% Section 16: Exact and closed $1$-forms: the first Stokes theorem} | Part: 3 | Depends on: see docs/REWRITE_PLAN.md | Status: todo | Original lines: 2538-2576
%!TEX root = ../../../main.tex
\subsection{Exact and closed $1$-forms: the first Stokes theorem}\label{sec:expl-closed}

For a function $f$ the differential is the $1$-form $\dd f=\partial_if\,\dd x^i$, so $\langle\dd f,v\rangle=v^i\partial_if=v(f)$ (the directional derivative). Along a curve:
\begin{align*}
\int_\gamma\dd f&=\int_a^b\partial_if(\gamma(t))\,\dot\gamma^i(t)\,\dd t &&\text{(definition \eqref{eq:lineint})}\\
&=\int_a^b\frac{\dd}{\dd t}f(\gamma(t))\,\dd t &&\text{(chain rule)}\\
&=f(\gamma(b))-f(\gamma(a)) &&\text{(fundamental theorem of calculus).}
\end{align*}
\begin{equation}\label{eq:ftc}
\int_\gamma\dd f=f(\gamma(b))-f(\gamma(a)).
\end{equation}
This is Stokes' theorem for $k=1$: the integral of $\dd f$ over the curve equals $f$ evaluated on the boundary $\partial\gamma=\{\gamma(b)\}-\{\gamma(a)\}$.

A form is \emph{exact} if $\omega=\dd f$ and \emph{closed} if $\dd\omega=0$.
\begin{itemize}
\item \emph{Exact $\Rightarrow$ closed.} $\dd(\dd f)=\dd(\partial_if\,\dd x^i)=\partial_j\partial_if\,\dd x^j\wedge\dd x^i=0$, because $\partial_j\partial_if$ is symmetric in $(i,j)$ while $\dd x^j\wedge\dd x^i$ is antisymmetric.
\item \emph{Exact $\Rightarrow$ integral zero on closed loops.} If $\gamma(b)=\gamma(a)$, \eqref{eq:ftc} gives $0$.
\item The converse ``closed $\Rightarrow$ exact'' fails in general: this is the starting point of de Rham cohomology (next example).
\end{itemize}

\subhead{How to compute $\dd$ of a $1$-form in the plane}
For $\omega=P\,\dd x+Q\,\dd y$:
\begin{align*}
\dd\omega&=\dd P\wedge\dd x+\dd Q\wedge\dd y\\
&=(P_x\dd x+P_y\dd y)\wedge\dd x+(Q_x\dd x+Q_y\dd y)\wedge\dd y\\
&=P_y\,\dd y\wedge\dd x+Q_x\,\dd x\wedge\dd y=(Q_x-P_y)\,\dd x\wedge\dd y,
\end{align*}
using $\dd x\wedge\dd x=\dd y\wedge\dd y=0$ and $\dd y\wedge\dd x=-\dd x\wedge\dd y$.

\begin{example}[The angle form]\label{ex:angle}
On $M=\R^2\setminus\{0\}$ let $\omega=\dfrac{-y\,\dd x+x\,\dd y}{x^2+y^2}=P\,\dd x+Q\,\dd y$, $P=\dfrac{-y}{x^2+y^2}$, $Q=\dfrac x{x^2+y^2}$.
\begin{description}[leftmargin=0pt,labelindent=0pt,font=\bfseries]
\item[Step 1: $\omega$ is closed.] With $r^2=x^2+y^2$: $Q_x=\dfrac{r^2-x\cdot2x}{r^4}=\dfrac{y^2-x^2}{r^4}$, $P_y=-\dfrac{r^2-y\cdot2y}{r^4}=\dfrac{y^2-x^2}{r^4}$, so $\dd\omega=(Q_x-P_y)\,\dd x\wedge\dd y=0$.
\item[Step 2: why it is called ``the angle''.] With $\theta=\arctan(y/x)$, $\partial_x\theta=\dfrac{-y/x^2}{1+y^2/x^2}=\dfrac{-y}{x^2+y^2}$ and $\partial_y\theta=\dfrac{1/x}{1+y^2/x^2}=\dfrac x{x^2+y^2}$, so $\dd\theta=\omega$ locally. (The function $\theta$ is multivalued around the origin.)
\item[Step 3: integrate on the unit circle] $\gamma(t)=(\cos t,\sin t)$, $0\le t\le2\pi$: velocity $\dot\gamma=(-\sin t,\cos t)$; on the circle $r=1$, $\omega_x=P=-\sin t$, $\omega_y=Q=\cos t$; contraction $\langle\omega_{\gamma(t)},\dot\gamma(t)\rangle=(-\sin t)(-\sin t)+(\cos t)(\cos t)=1$; integral $\oint_\gamma\omega=\int_0^{2\pi}1\,\dd t=2\pi$.
\item[Step 4: conclusion.] If $\omega=\dd f$ with $f$ single-valued on $M$, then \eqref{eq:ftc} would give $\oint_\gamma\omega=f(\gamma(2\pi))-f(\gamma(0))=0$, contradicting Step 3. Hence $\omega$ is closed but not exact on $M$: there is no global angle function. The number $\frac1{2\pi}\oint\omega$ is the \emph{winding number}, and $H^1_{\mathrm{dR}}(\R^2\setminus\{0\})\cong\R$.
\end{description}
\end{example}

% Section 16.6: A local meaning of $\dd\omega$: circulation per unit area | Part: 3 | Depends on: sec:fields, sec:tangent | Status: todo | Original lines: 2577-2608
%!TEX root = ../../../main.tex
\subsection{A local meaning of $\dd\omega$: circulation per unit area}\label{sec:circulation}

Let us compute the line integral of $\omega=\omega_1\dd x^1+\omega_2\dd x^2$ around a tiny rectangle with lower-left corner $(x^1,x^2)$ and sides $\varepsilon$ (direction 1) and $\delta$ (direction 2), counter-clockwise.

\Step{1}{the four edges} On each short edge, $\int\omega\approx\omega(\text{midpoint of the edge})\cdot(\text{displacement})$:
\begin{align*}
\text{bottom }(x^1,x^2)\to(x^1{+}\varepsilon,x^2)&:\ \omega_1\bigl(x^1{+}\tfrac\varepsilon2,x^2\bigr)\varepsilon,\\
\text{right }(x^1{+}\varepsilon,x^2)\to(x^1{+}\varepsilon,x^2{+}\delta)&:\ \omega_2\bigl(x^1{+}\varepsilon,x^2{+}\tfrac\delta2\bigr)\delta,\\
\text{top }(x^1{+}\varepsilon,x^2{+}\delta)\to(x^1,x^2{+}\delta)&:\ -\omega_1\bigl(x^1{+}\tfrac\varepsilon2,x^2{+}\delta\bigr)\varepsilon,\\
\text{left }(x^1,x^2{+}\delta)\to(x^1,x^2)&:\ -\omega_2\bigl(x^1,x^2{+}\tfrac\delta2\bigr)\delta.
\end{align*}
\Step{2}{group the terms}
\[
\oint\omega\approx\varepsilon\Bigl[\omega_1\bigl(x^1{+}\tfrac\varepsilon2,x^2\bigr)-\omega_1\bigl(x^1{+}\tfrac\varepsilon2,x^2{+}\delta\bigr)\Bigr]+\delta\Bigl[\omega_2\bigl(x^1{+}\varepsilon,x^2{+}\tfrac\delta2\bigr)-\omega_2\bigl(x^1,x^2{+}\tfrac\delta2\bigr)\Bigr].
\]
\Step{3}{Taylor expansion} $\omega_1(x^1{+}\frac\varepsilon2,x^2{+}\delta)-\omega_1(x^1{+}\frac\varepsilon2,x^2)\approx\delta\,\partial_2\omega_1$ and $\omega_2(x^1{+}\varepsilon,x^2{+}\frac\delta2)-\omega_2(x^1,x^2{+}\frac\delta2)\approx\varepsilon\,\partial_1\omega_2$. Therefore, up to terms of higher order than $\varepsilon\delta$,
\[
\oint\omega\approx-\varepsilon\delta\,\partial_2\omega_1+\varepsilon\delta\,\partial_1\omega_2=\varepsilon\delta\,\bigl(\partial_1\omega_2-\partial_2\omega_1\bigr).
\]
(Evaluating at the starting point of each edge would leave extra terms $-\varepsilon^2\partial_1\omega_1-\delta^2\partial_2\omega_2$, which is why the midpoint is used.)
\Step{4}{recognise $\dd\omega$} By the plane formula above, $\dd\omega=(\partial_1\omega_2-\partial_2\omega_1)\,\dd x^1\wedge\dd x^2$, and $(\dd x^1\wedge\dd x^2)(\varepsilon\partial_1,\delta\partial_2)=\varepsilon\delta$. Hence
\begin{equation}\label{eq:circ}
\oint_{\partial(\text{rectangle})}\omega\approx\dd\omega(\varepsilon\partial_1,\delta\partial_2).
\end{equation}
For a parallelogram spanned by $\varepsilon u,\delta v$ with arbitrary vector fields, the same computation gives $\varepsilon\delta\,\dd\omega(u,v)$ with $\dd\omega(u,v)=u(\omega(v))-v(\omega(u))-\omega([u,v])$ (the invariant formula of Theorem~\ref{thm:dprops}); for $u=\partial_1$, $v=\partial_2$ the bracket vanishes and we recover $\partial_1\omega_2-\partial_2\omega_1$.
\Step{5}{gluing (the idea of Stokes' theorem)} Chop a big region into such small rectangles. Every interior edge is traversed twice, once in each direction by the two neighbouring rectangles, so the contributions cancel in pairs; only the outer boundary survives:
\[
\oint_{\partial(\text{region})}\omega=\sum_{\text{rectangles}}\oint\omega\approx\sum\dd\omega(\dots)\ \longrightarrow\ \int_{\text{region}}\dd\omega .
\]
Thus $\dd\omega$ is the ``circulation per unit area'': it measures how far the sheets of $\omega$ fail to close up.


% Section 17: From $1$-forms to $k$-forms: integration is full contraction} | Part: 3 | Depends on: see docs/REWRITE_PLAN.md | Status: todo | Original lines: 2609-2612
%!TEX root = ../../../main.tex
% =====================================================================
\section{From $1$-forms to $k$-forms: integration is full contraction}\label{sec:intk}
% =====================================================================
% Section 17: 01-full-contraction | Part: 3 | Depends on: see docs/REWRITE_PLAN.md | Status: todo | Original lines: 2613-2635
%!TEX root = ../../../main.tex
Let $\phi:U\subset\R^k\to M$ be a parametrised $k$-dimensional piece of $M$, and $\alpha$ a $k$-form. Then
\begin{equation}\label{eq:intk}
\int_\phi\alpha:=\int_U\phi\pb\alpha=\int_U\alpha_{\phi(u)}\bigl(\partial_1\phi,\dots,\partial_k\phi\bigr)\,\dd u^1\cdots\dd u^k,
\end{equation}
that is, evaluate $\alpha$ on the tangent $k$-frame of the parametrisation (a \emph{full contraction}) and integrate the resulting function with the Lebesgue measure. For $k=1$ this is \eqref{eq:lineint}.
The same reasoning as in Proposition~\ref{prop:reparam} shows invariance under orientation-preserving reparametrisations (the Jacobian determinant appears through the multilinear \emph{alternating} character of $\alpha$), which is why the integrands of oriented $k$-dimensional objects must be alternating $k$-forms. On an oriented manifold one globalises with a partition of unity.

\emph{Why a determinant appears.} For a $2$-form $\alpha=a\,\dd x\wedge\dd y$ and a parametrisation $\phi(u,v)=(x(u,v),y(u,v))$, the tangent frame is $\partial_u\phi=(x_u,y_u)$, $\partial_v\phi=(x_v,y_v)$, so
\[
\alpha(\partial_u\phi,\partial_v\phi)=a\,(x_uy_v-y_ux_v)=a\det D\phi,
\]
where $D\phi$ is the Jacobian matrix. Therefore \eqref{eq:intk} is the change-of-variables formula $\int_\phi a\,\dd x\wedge\dd y=\iint a(\phi(u,v))\det D\phi\,\dd u\,\dd v$ with the \emph{sign} of the Jacobian kept: orientation-preserving reparametrisations ($\det>0$) leave it unchanged and orientation-reversing ones flip its sign.

\begin{example}[A surface integral, step by step]\label{ex:paraboloid}
Integrate $\alpha=z\,\dd x\wedge\dd y$ over the paraboloid cap $\phi(u,v)=(u,v,u^2+v^2)$, $u^2+v^2\le1$.
\begin{enumerate}
\item Pull back the coordinate functions: $\phi\pb x=u$, $\phi\pb y=v$, $\phi\pb z=u^2+v^2$; since pull-back commutes with $\dd$ and $\wedge$, $\phi\pb\dd x=\dd u$ and $\phi\pb\dd y=\dd v$.
\item Hence $\phi\pb\alpha=(u^2+v^2)\,\dd u\wedge\dd v$.
\item Evaluate on the frame $(\partial_u,\partial_v)$: $(\dd u\wedge\dd v)(\partial_u,\partial_v)=1$, so $\int_\phi\alpha=\iint_{u^2+v^2\le1}(u^2+v^2)\,\dd u\,\dd v$.
\item Polar coordinates (the factor $r$ is the determinant of Example~\ref{ex:polarwedge}): $\displaystyle\int_0^{2\pi}\!\!\int_0^1r^2\cdot r\,\dd r\,\dd\theta=2\pi\cdot\tfrac14=\tfrac\pi2$.
\end{enumerate}
\end{example}

% Section 17: Partial contraction and flux} | Part: 3 | Depends on: see docs/REWRITE_PLAN.md | Status: todo | Original lines: 2636-2662
%!TEX root = ../../../main.tex
\subsection{Partial contraction and flux}\label{sec:flux}

Partial contractions appear whenever a vector field is to be integrated over a hypersurface. Let $\mu$ be an $n$-form (\eg a volume form) and $X$ a vector field. Then $\iota_X\mu$ is an $(n-1)$-form, and for tangent vectors $w_2,\dots,w_n$ to a hypersurface $\Sigma$,
\[
(\iota_X\mu)(w_2,\dots,w_n)=\mu(X,w_2,\dots,w_n)
\]
is the signed volume of the parallelepiped spanned by $X,w_2,\dots,w_n$: the volume of material carried by the flow of $X$ through the element of area spanned by $w_2,\dots,w_n$. Hence
\begin{equation}\label{eq:flux}
\text{flux of }X\text{ through }\Sigma=\int_\Sigma\iota_X\mu .
\end{equation}
\emph{No normal vector and no metric have been used;} with a Riemannian metric this will turn out to equal $\int_\Sigma\langle X,N\rangle\,\dd A$ (Lemma~\ref{lem:fluxbdry}).

\emph{Computing $\iota_X\mu$ in $\R^3$ (steps).} Let $\mu=\dd x\wedge\dd y\wedge\dd z$ and $X=X^1\partial_x+X^2\partial_y+X^3\partial_z$. Using the antiderivation rule on $\dd x\wedge(\dd y\wedge\dd z)$:
\begin{enumerate}
\item $\iota_X(\dd y\wedge\dd z)=X^2\dd z-X^3\dd y$ (same computation as in Example~\ref{ex:interiorplane}).
\item $\iota_X\mu=(\iota_X\dd x)\,\dd y\wedge\dd z-\dd x\wedge\iota_X(\dd y\wedge\dd z)=X^1\dd y\wedge\dd z-\dd x\wedge(X^2\dd z-X^3\dd y)$.
\item Expanding and using $-\dd x\wedge\dd z=\dd z\wedge\dd x$:
\begin{equation}\label{eq:iotavol}
\iota_X\mu=X^1\,\dd y\wedge\dd z+X^2\,\dd z\wedge\dd x+X^3\,\dd x\wedge\dd y .
\end{equation}
\end{enumerate}
\emph{Check:} on a horizontal plane $z=\mathrm{const}$ oriented by $(\partial_x,\partial_y)$, $\dd z=0$ so only $X^3\dd x\wedge\dd y$ survives: the flux is $\iint X^3\,\dd x\,\dd y$, the component of $X$ along the normal $\partial_z$, as expected.

\begin{synthesis}[the same formula in geometric algebra]
Under $\dd x^i\leftrightarrow e_i$, \eqref{eq:iotavol} is the bivector $X^1e_2e_3+X^2e_3e_1+X^3e_1e_2=X\pss$ (Proposition~\ref{prop:threelang}): the flux $2$-form of $X$ is the \emph{dual} of the vector $X$. This is the geometric content of ``a vector in $\R^3$ is secretly a plane element'' (Section~\ref{sec:duality}).
\end{synthesis}


% Section 18: Stokes' theorem as it is understood today | Part: 3 | Depends on: sec:intk, sec:fields | Status: todo | Original lines: 2663-2666
%!TEX root = ../../../main.tex
% =====================================================================
\section{Stokes' theorem as it is understood today}\label{sec:stokes}
% =====================================================================
% Section 18: Statement} | Part: 3 | Depends on: see docs/REWRITE_PLAN.md | Status: todo | Original lines: 2667-2680
%!TEX root = ../../../main.tex
\subsection{Statement}

Let $M$ be an oriented smooth $n$-manifold with boundary, and give $\partial M$ the \emph{induced orientation} (outward-pointing vector first: a basis $(w_2,\dots,w_n)$ of $T_p\partial M$ is positive iff $(N,w_2,\dots,w_n)$ is positive in $T_pM$). Let $j:\partial M\hookrightarrow M$ be the inclusion (the letter $\iota$ is reserved for the interior product).

\begin{theorem}[Generalised Stokes theorem]\label{thm:stokes}
For every compactly supported smooth $(n-1)$-form $\omega$ on $M$,
\begin{equation}\label{eq:stokes}
\int_M\dd\omega=\int_{\partial M}j\pb\omega .
\end{equation}
(If $\partial M=\emptyset$, the right-hand side is zero.)
\end{theorem}

The theorem contains, in a single line, the fundamental theorem of calculus ($n=1$), Green's theorem ($n=2$), the Kelvin--Stokes theorem ($n=2$ surface in $\R^3$) and the divergence theorem ($n=3$, or any $n$).

% Section 18: Warm-up: Stokes' theorem on a rectangle (Green's theorem)} | Part: 3 | Depends on: see docs/REWRITE_PLAN.md | Status: todo | Original lines: 2681-2704
%!TEX root = ../../../main.tex
\subsection{Warm-up: Stokes' theorem on a rectangle (Green's theorem)}\label{sec:green}

Before the general proof, let us see the mechanism in the simplest case. Let $\omega=P\,\dd x+Q\,\dd y$ on the rectangle $R=[a,b]\times[c,d]$ with counter-clockwise boundary.

\Step{1}{left-hand side} By the plane formula of Section~\ref{sec:expl-closed}, $\dd\omega=(Q_x-P_y)\,\dd x\wedge\dd y$, hence $\int_R\dd\omega=\iint_R(Q_x-P_y)\,\dd x\,\dd y$.

\Step{2}{fundamental theorem of calculus in each variable}
\[
\iint_RQ_x\,\dd x\,\dd y=\int_c^d\bigl[Q(b,y)-Q(a,y)\bigr]\dd y,\qquad
\iint_RP_y\,\dd x\,\dd y=\int_a^b\bigl[P(x,d)-P(x,c)\bigr]\dd x .
\]

\Step{3}{right-hand side: the four edges}
\begin{itemize}
\item bottom ($y=c$, $x:a\to b$, $\dd y=0$): $\int_a^bP(x,c)\,\dd x$; right ($x=b$, $y:c\to d$, $\dd x=0$): $\int_c^dQ(b,y)\,\dd y$;
\item top ($y=d$, $x:b\to a$): $-\int_a^bP(x,d)\,\dd x$; left ($x=a$, $y:d\to c$): $-\int_c^dQ(a,y)\,\dd y$.
\end{itemize}

\Step{4}{compare} The sum of the four edges is $\int_c^d[Q(b,y)-Q(a,y)]\dd y-\int_a^b[P(x,d)-P(x,c)]\dd x$, which is exactly Step 2. Hence $\oint_{\partial R}\omega=\int_R\dd\omega$.

\numcheck $\omega=-y\,\dd x+x\,\dd y$ on $[0,1]^2$: $Q_x-P_y=1-(-1)=2$, so $\int_R\dd\omega=2$. Boundary: bottom $y=0$: $0$; right $x=1$: $\int_0^11\,\dd y=1$; top $y=1$: $-\int_0^1(-1)\,\dd x=1$; left $x=0$: $0$. Total $2$ \checkmark.

The general proof does the same thing in $n$ variables: the fundamental theorem of calculus in one variable, and sign bookkeeping.

% Section 18.3: The proof of Theorem~\ref{thm:stokes}, step by step | Part: 3 | Depends on: sec:intk, sec:fields | Status: todo | Original lines: 2705-2754
%!TEX root = ../../../main.tex
\subsection{The proof of Theorem~\ref{thm:stokes}, step by step}\label{sec:stokesproof}

\Step{0}{reduction to a chart} Take a partition of unity $\{\rho_\alpha\}$ subordinate to a cover by oriented charts and write $\omega=\sum_\alpha\rho_\alpha\omega$. Both sides of \eqref{eq:stokes} are linear in $\omega$ and invariant under orientation-preserving diffeomorphisms (naturality \eqref{eq:natural}, extended to $k$-forms), so it suffices to prove it for forms supported in one chart, \ie on $\R^n$ or on the half-space $\mathbb H^n=\{x^n\ge0\}$. For $\R^n$ there is no boundary and the claim is $\int\dd\omega=0$. This follows from the computation below, where now every term $i=1,\dots,n$ vanishes: each integral $\int_{-\infty}^\infty\partial_i\omega_i\,\dd x^i$ is zero because $\omega_i$ has compact support.

\Step{1}{write $\omega$} Any $(n-1)$-form on $\mathbb H^n$ is $\omega=\sum_{i=1}^n\omega_i\,\dd x^1\wedge\cdots\wedge\widehat{\dd x^i}\wedge\cdots\wedge\dd x^n$ (the hat means ``omit''), with $\omega_i$ smooth of compact support.

\Step{2}{compute $\dd\omega$}
\[
\dd\Bigl(\omega_i\,\dd x^1\wedge\cdots\wedge\widehat{\dd x^i}\wedge\cdots\wedge\dd x^n\Bigr)=\sum_j\partial_j\omega_i\,\dd x^j\wedge\dd x^1\wedge\cdots\wedge\widehat{\dd x^i}\wedge\cdots\wedge\dd x^n .
\]
If $j\ne i$ the factor $\dd x^j$ already appears in the product, so the term vanishes ($\dd x^j\wedge\dd x^j=0$). For $j=i$, moving $\dd x^i$ past the $i-1$ factors $\dd x^1,\dots,\dd x^{i-1}$ gives the sign $(-1)^{i-1}$. Hence $\dd\omega=\sum_i(-1)^{i-1}\partial_i\omega_i\,\dd x^1\wedge\cdots\wedge\dd x^n$.

\Step{3}{integrate over $\mathbb H^n$; Fubini and the fundamental theorem of calculus}
\begin{itemize}
\item For $i<n$, integrate first in $x^i\in\R$: $\int_{-\infty}^{\infty}\partial_i\omega_i\,\dd x^i=\omega_i\big|_{-\infty}^{+\infty}=0$ because $\omega_i$ has compact support. These terms contribute $0$.
\item For $i=n$, integrate in $x^n\in[0,\infty)$: $\int_0^\infty\partial_n\omega_n\,\dd x^n=0-\omega_n(x',0)=-\omega_n(x',0)$, where $x'=(x^1,\dots,x^{n-1})$.
\end{itemize}
Therefore, with the sign $(-1)^{n-1}$ of the $i=n$ term, $\int_{\mathbb H^n}\dd\omega=(-1)^{n-1}\cdot(-1)\int_{\R^{n-1}}\omega_n(x',0)\,\dd x'=(-1)^n\int_{\R^{n-1}}\omega_n(x',0)\,\dd x'$.

\Step{4}{the boundary side} On $\partial\mathbb H^n=\{x^n=0\}$ we have $\dd x^n\mapsto0$ under pull-back, so every term of $\omega$ containing $\dd x^n$ (all $i\neq n$) vanishes, and $j\pb\omega=\omega_n(x',0)\,\dd x^1\wedge\cdots\wedge\dd x^{n-1}$.

\Step{5}{orientation of $\partial\mathbb H^n$} The outward normal is $-\partial_n$. A basis $(w_1,\dots,w_{n-1})$ of $T\partial\mathbb H^n$ is positive iff $(-\partial_n,w_1,\dots,w_{n-1})$ is positive in $\R^n$. For the standard basis $w_j=\partial_j$ we get
$\det(-e_n,e_1,\dots,e_{n-1})=-\det(e_n,e_1,\dots,e_{n-1})=-(-1)^{n-1}=(-1)^n$, because moving $e_n$ from the first to the last position is a cyclic permutation of $n$ vectors, of sign $(-1)^{n-1}$. So the induced orientation on $\partial\mathbb H^n\cong\R^{n-1}$ is $(-1)^n$ times the standard one, and $\int_{\partial\mathbb H^n}j\pb\omega=(-1)^n\int_{\R^{n-1}}\omega_n(x',0)\,\dd x'$, which equals the result of Step 3. \hfill$\blacksquare$

\begin{specialcase}[sanity checks of the sign: $n=1$ and $n=2$]
$n=1$: $\mathbb H^1=[0,\infty)$, the boundary is the point $0$ with sign $(-1)^1=-1$, and indeed $\int_0^\infty f'\,\dd x=-f(0)$.
$n=2$: the sign is $+1$, the upper half-plane is traversed in the $+x$ direction (region on the left), consistent with the counter-clockwise convention of Green's theorem.
\end{specialcase}

% Figure | used in parts/03-integration/03-stokes/03-proof.tex
%!TEX root = ../../main.tex
\begin{figure}[ht]
\centering
\begin{tikzpicture}[>=Stealth,scale=1.05]
  \fill[blue!12] (0,0) circle (1.9);
  \fill[white] (0,0) circle (0.65);
  \draw[navy] (0,0) circle (1.9);
  \draw[navy] (0,0) circle (0.65);
  \node[navy] at (-1.15,0.55) {$M$};
  % outer boundary: ccw arrow
  \draw[->,navy,thick] (60:1.9) arc (60:75:1.9);
  \draw[->,red,thick] (45:1.9)--(45:2.75) node[above right] {$N$};
  % inner boundary: clockwise arrow, N points into the hole
  \draw[->,navy,thick] (200:0.65) arc (200:185:0.65);
  \draw[->,red,thick] (330:0.65)--(330:0.15);
\end{tikzpicture}

\small $N=$ outward normal of $M$ (on the inner circle it points into the hole)
\caption{An annulus $M$ with its induced boundary orientation: the outer circle is traversed counter-clockwise and the inner circle clockwise, so that the region lies to the left. Applied to the angle form of Example~\ref{ex:angle} ($\dd\omega=0$) Stokes gives $\oint_{\mathrm{outer}}\omega-\oint_{\mathrm{inner\ (ccw)}}\omega=0$: both equal $2\pi$.}
\label{fig:annulus}
\end{figure}


% Section 18.4: The structural meaning: $\dd$ is the adjoint of $\partial$ | Part: 3 | Depends on: sec:intk, sec:fields | Status: todo | Original lines: 2755-2768
%!TEX root = ../../../main.tex
\subsection{The structural meaning: $\dd$ is the adjoint of $\partial$}\label{sec:adjoint}

The modern reading of \eqref{eq:stokes} is obtained by introducing the pairing between $k$-forms and $k$-chains $c$ (formal sums of oriented $k$-dimensional pieces): $\langle c,\alpha\rangle:=\int_c\alpha$. Then Stokes' theorem says
\begin{equation}\label{eq:adjoint}
\langle c,\dd\alpha\rangle=\langle\partial c,\alpha\rangle,
\end{equation}
\ie the exterior derivative and the boundary operator are \emph{adjoint} to each other. Consequences:
\begin{itemize}
\item $\partial^2=0$ (the boundary of a boundary is empty) is the dual of $\dd^2=0$, and each explains the other. Indeed, applying \eqref{eq:adjoint} twice: $\langle c,\dd\dd\alpha\rangle=\langle\partial c,\dd\alpha\rangle=\langle\partial\partial c,\alpha\rangle$. If $\partial\partial c=0$ the right-hand side vanishes for every $c$, and a form whose integral over every small piece is zero must be zero: $\dd\dd\alpha=0$.
\item Integration is a chain map from the de Rham complex $(\Omega^\bullet(M),\dd)$ to the singular cochain complex, hence it descends to cohomology. de Rham's theorem (1931): this map $H^k_{\mathrm{dR}}(M)\to H^k(M;\R)$ is an isomorphism. Closed forms integrate to the same value over homologous cycles (``periods''), as in Example~\ref{ex:angle}.
\item Stokes' theorem is \emph{metric-free} and uses only a smooth structure and an orientation. Metrics enter only when we translate forms into vector fields.
\item Applications reach well beyond calculus: degree theory and the Brouwer fixed-point theorem, the Poincar\'e lemma and Poincar\'e duality, conservation laws, and Maxwell's equations $\dd F=0$, $\dd\hodge F=J$, whose integral forms are the laws of Faraday and Gauss.
\end{itemize}

% Section 18.5: Classical corollaries | Part: 3 | Depends on: sec:intk, sec:fields | Status: todo | Original lines: 2769-2819
%!TEX root = ../../../main.tex
\subsection{Classical corollaries}\label{sec:classical}

Let $\R^3$ carry its Euclidean metric $g$, orientation and volume form $\vol=\dd x\wedge\dd y\wedge\dd z$. With the musical isomorphisms $\flat,\sharp$ and the Hodge star $\hodge$ one has the dictionary
\[
\operatorname{grad}f=(\dd f)^\sharp,\qquad
\operatorname{curl}F=\bigl(\hodge\,\dd F^\flat\bigr)^\sharp,\qquad
\operatorname{div}F=\hodge\,\dd\hodge F^\flat,\qquad
\iota_F\vol=\hodge F^\flat .
\]
Applying Theorem~\ref{thm:stokes} to forms of degree $0,1,2$ yields:

\begin{center}
\small
\renewcommand{\arraystretch}{1.4}
\begin{tabular}{@{}p{2.5cm}p{4.5cm}p{7.0cm}@{}}
\toprule
\textbf{Form of Stokes} & \textbf{Differential forms} & \textbf{Classical statement}\\
\midrule
$n=1$ & $\int_a^b\dd f=f(b)-f(a)$ & fundamental theorem of calculus / gradient theorem\\
$n=2$ (planar $D$) & $\int_D\dd\omega=\oint_{\partial D}\omega$ & Green: $\oint P\,\dd x+Q\,\dd y=\iint(Q_x-P_y)\,\dd A$\\
$n=2$ (surface $S\subset\R^3$) & $\int_S\dd F^\flat=\oint_{\partial S}F^\flat$ & Kelvin--Stokes: $\oint F\cdot\dd r=\iint(\operatorname{curl}F)\cdot n\,\dd A$\\
$n=3$ & $\int_\Omega\dd(\iota_F\vol)=\int_{\partial\Omega}\iota_F\vol$ & Gauss: $\iiint\operatorname{div}F\,\dd V=\iint_{\partial\Omega}F\cdot n\,\dd A$\\
\bottomrule
\end{tabular}
\end{center}

\subhead{Verification of the dictionary, step by step}
In $\R^3$ with the Euclidean metric, lowering an index changes nothing numerically: $F^\flat=F_x\dd x+F_y\dd y+F_z\dd z$. Recall \eqref{eq:iotavol}: $\iota_G\vol=G_x\,\dd y\wedge\dd z+G_y\,\dd z\wedge\dd x+G_z\,\dd x\wedge\dd y$.

\emph{(i) Curl.} Compute $\dd F^\flat$ term by term, using $\dd x\wedge\dd x=0$:
\begin{align*}
\dd(F_x\dd x)&=\partial_yF_x\,\dd y\wedge\dd x+\partial_zF_x\,\dd z\wedge\dd x=-\partial_yF_x\,\dd x\wedge\dd y+\partial_zF_x\,\dd z\wedge\dd x,\\
\dd(F_y\dd y)&=\partial_xF_y\,\dd x\wedge\dd y+\partial_zF_y\,\dd z\wedge\dd y=\partial_xF_y\,\dd x\wedge\dd y-\partial_zF_y\,\dd y\wedge\dd z,\\
\dd(F_z\dd z)&=\partial_xF_z\,\dd x\wedge\dd z+\partial_yF_z\,\dd y\wedge\dd z=-\partial_xF_z\,\dd z\wedge\dd x+\partial_yF_z\,\dd y\wedge\dd z.
\end{align*}
Adding and collecting the three basis $2$-forms:
\[
\dd F^\flat=(\partial_yF_z-\partial_zF_y)\,\dd y\wedge\dd z+(\partial_zF_x-\partial_xF_z)\,\dd z\wedge\dd x+(\partial_xF_y-\partial_yF_x)\,\dd x\wedge\dd y=\iota_{\operatorname{curl}F}\vol,
\]
since the three coefficients are exactly the components of $\operatorname{curl}F$. Because $\hodge\hodge=1$ in $\R^3$ and $\hodge(\dd y\wedge\dd z)=\dd x$, etc., this is the same as $\operatorname{curl}F=(\hodge\,\dd F^\flat)^\sharp$.

\emph{(ii) Divergence.} $\iota_F\vol=F_x\dd y\wedge\dd z+F_y\dd z\wedge\dd x+F_z\dd x\wedge\dd y$. Differentiating, only the derivative with respect to the variable missing from each $2$-form survives:
$\dd(F_x\dd y\wedge\dd z)=\partial_xF_x\,\dd x\wedge\dd y\wedge\dd z$, $\dd(F_y\dd z\wedge\dd x)=\partial_yF_y\,\dd y\wedge\dd z\wedge\dd x$, $\dd(F_z\dd x\wedge\dd y)=\partial_zF_z\,\dd z\wedge\dd x\wedge\dd y$,
and the three $3$-forms equal $\dd x\wedge\dd y\wedge\dd z$ (cyclic permutations of three $1$-forms are even). So $\dd(\iota_F\vol)=(\operatorname{div}F)\,\vol$.

\emph{(iii) Kelvin--Stokes.} For an oriented surface $S\subset\R^3$ with unit normal $n$ the restriction satisfies $\iota_G\vol|_S=\langle G,n\rangle\,\dd A$ (same computation as in Lemma~\ref{lem:fluxbdry} below, with $e_1=n$ and $e_2,e_3$ tangent to $S$). Then
\[
\oint_{\partial S}F\cdot\dd r=\oint_{\partial S}F^\flat\overset{\text{Stokes}}=\int_S\dd F^\flat=\int_S\iota_{\operatorname{curl}F}\vol=\iint_S\langle\operatorname{curl}F,n\rangle\,\dd A .
\]
\emph{(iv) Gauss.} $\displaystyle\int_\Omega(\operatorname{div}F)\,\dd V\overset{\text{(ii)}}=\int_\Omega\dd(\iota_F\vol)\overset{\text{Stokes}}=\int_{\partial\Omega}\iota_F\vol=\iint_{\partial\Omega}\langle F,n\rangle\,\dd A$, again by the restriction formula.


% Section 19: The Gauss (divergence) theorem | Part: 3 | Depends on: sec:stokes, sec:interior, sec:metric | Status: todo | Original lines: 2820-2823
%!TEX root = ../../../main.tex
% =====================================================================
\section{The Gauss (divergence) theorem}\label{sec:gauss}
% =====================================================================
% Section 19: Divergence is defined by a volume form, not by a metric} | Part: 3 | Depends on: see docs/REWRITE_PLAN.md | Status: todo | Original lines: 2824-2845
%!TEX root = ../../../main.tex
\subsection{Divergence is defined by a volume form, not by a metric}\label{sec:divmu}

Let $\mu$ be a nowhere-vanishing $n$-form on $M$ (a \emph{volume form}; in the Riemannian case $\mu=\vol_g=\sqrt{|g|}\,\dd x^1\wedge\cdots\wedge\dd x^n$). The \emph{divergence} of a vector field $X$ is defined by
\begin{equation}\label{eq:Lmu}
\Lie_X\mu=(\operatorname{div}_\mu X)\,\mu,
\end{equation}
that is, as the relative rate of change of volume under the flow of $X$. Since $\dd\mu=0$ (it is an $(n+1)$-form), Cartan's formula \eqref{eq:cartan} gives
\begin{equation}\label{eq:divmu}
\dd(\iota_X\mu)=\Lie_X\mu=(\operatorname{div}_\mu X)\,\mu .
\end{equation}
\emph{Coordinate formula, step by step.} Let $\mu=\sqrt{|g|}\,\dd x^1\wedge\cdots\wedge\dd x^n$ and $X=X^i\partial_i$.
\begin{enumerate}
\item Contract (antiderivation rule applied repeatedly, as in Section~\ref{sec:flux}): \[\iota_X\mu=\sqrt{|g|}\sum_i(-1)^{i-1}X^i\,\dd x^1\wedge\cdots\wedge\widehat{\dd x^i}\wedge\cdots\wedge\dd x^n.\]
\item Differentiate: in $\dd(\sqrt{|g|}X^i)\wedge\dd x^1\wedge\cdots\wedge\widehat{\dd x^i}\wedge\cdots$ only $\partial_i(\sqrt{|g|}X^i)\,\dd x^i$ survives (the other $\dd x^j$ are repeated), and moving $\dd x^i$ to its place gives another $(-1)^{i-1}$. The two signs cancel: $\dd(\iota_X\mu)=\sum_i\partial_i(\sqrt{|g|}X^i)\,\dd x^1\wedge\cdots\wedge\dd x^n$.
\item Divide by $\mu$:
\begin{equation}\label{eq:divcoord}
\operatorname{div}_\mu X=\frac1{\sqrt{|g|}}\,\partial_i\bigl(\sqrt{|g|}\,X^i\bigr),
\end{equation}
which reduces to $\partial_iX^i$ in Cartesian coordinates ($|g|=1$).
\end{enumerate}
So the divergence theorem is nothing but Theorem~\ref{thm:stokes} applied to the $(n-1)$-form $\omega=\iota_X\mu$ obtained from $X$ by (partial) contraction.

% Section 19: The boundary term: from $\iota_X\mu$ to $\langle X,N\rangle\,\dd A$} | Part: 3 | Depends on: see docs/REWRITE_PLAN.md | Status: todo | Original lines: 2846-2864
%!TEX root = ../../../main.tex
\subsection{The boundary term: from $\iota_X\mu$ to $\langle X,N\rangle\,\dd A$}\label{sec:boundaryterm}

\begin{lemma}\label{lem:fluxbdry}
Let $(M,g)$ be an oriented Riemannian manifold with boundary, $N$ the outward unit normal and $\dd A$ the Riemannian area form of $\partial M$ (with the induced orientation). Then
\[
j\pb\bigl(\iota_X\vol_g\bigr)=\langle X,N\rangle\,\dd A .
\]
\end{lemma}
\begin{proof}
\Step{1}{adapted frame} At a boundary point choose an oriented orthonormal frame $e_1=N,e_2,\dots,e_n$ ($e_2,\dots,e_n$ tangent to $\partial M$) with dual coframe $e^1,\dots,e^n$. Then $\vol_g=e^1\wedge\cdots\wedge e^n$. By the induced-orientation convention $(e_2,\dots,e_n)$ is a positive orthonormal basis of $T\partial M$, so $\dd A=j\pb(e^2\wedge\cdots\wedge e^n)$.
\Step{2}{contract} Apply $\iota_X$ one factor at a time: $\iota_X(e^i\wedge\beta)=X^i\beta-e^i\wedge\iota_X\beta$. For $n=3$:
\begin{align*}
\iota_X(e^1\wedge e^2\wedge e^3)&=X^1e^2\wedge e^3-e^1\wedge\iota_X(e^2\wedge e^3)\\
&=X^1e^2\wedge e^3-e^1\wedge(X^2e^3-X^3e^2)=X^1e^2\wedge e^3-X^2e^1\wedge e^3+X^3e^1\wedge e^2 .
\end{align*}
In general, $\iota_X\vol_g=\sum_{i=1}^n(-1)^{i-1}X^i\,e^1\wedge\cdots\wedge\widehat{e^i}\wedge\cdots\wedge e^n$.
\Step{3}{restrict to $\partial M$} The covector $e^1=\langle N,\cdot\rangle$ vanishes on all tangent vectors $e_2,\dots,e_n$, so $j\pb e^1=0$ and every term containing $e^1$ (all $i\ne1$) restricts to zero. Only $i=1$ survives: $j\pb(\iota_X\vol_g)=X^1j\pb(e^2\wedge\cdots\wedge e^n)=X^1\dd A$, and $X^1=\langle X,e_1\rangle=\langle X,N\rangle$. In the $n=3$ example above this reads $j\pb(\cdots)=X^1e^2\wedge e^3$.
\end{proof}

% Section 19: Statement and proof} | Part: 3 | Depends on: see docs/REWRITE_PLAN.md | Status: todo | Original lines: 2865-2891
%!TEX root = ../../../main.tex
\subsection{Statement and proof}

\begin{theorem}[Gauss--Ostrogradsky, divergence theorem]\label{thm:gauss}
Let $(M,g)$ be a compact oriented Riemannian $n$-manifold with boundary and $X$ a smooth vector field. Then
\[
\int_M\operatorname{div}X\,\dd V=\int_{\partial M}\langle X,N\rangle\,\dd A .
\]
\end{theorem}
\begin{proof}
By \eqref{eq:divmu}, $\operatorname{div}X\,\vol_g=\dd(\iota_X\vol_g)$. Apply Theorem~\ref{thm:stokes} to $\omega=\iota_X\vol_g$ and use Lemma~\ref{lem:fluxbdry}.
\end{proof}

\begin{example}[The position field on the unit ball]\label{ex:ball}
Let $X=x\partial_x+y\partial_y+z\partial_z$ on the unit ball $B\subset\R^3$.
\begin{enumerate}
\item \emph{Volume side.} $\operatorname{div}X=\partial_xx+\partial_yy+\partial_zz=3$, so $\int_B3\,\dd V=3\cdot\frac{4\pi}3=4\pi$.
\item \emph{Boundary side, with the formula of the lemma.} On the unit sphere the outward normal is $N=(x,y,z)=X$, hence $\langle X,N\rangle=x^2+y^2+z^2=1$, and the flux is $\int_{S^2}1\,\dd A=4\pi$.
\item \emph{Boundary side, directly with forms.} $\iota_X\vol=x\,\dd y\wedge\dd z+y\,\dd z\wedge\dd x+z\,\dd x\wedge\dd y$ (formula \eqref{eq:iotavol}). At the north pole $(0,0,1)$ the tangent plane is spanned by $\partial_x,\partial_y$, where $\dd z=0$, so the restriction is $z\,\dd x\wedge\dd y=\dd x\wedge\dd y=\dd A$, again with coefficient $1$.
\end{enumerate}
Both sides equal $4\pi$ \checkmark.
\end{example}

\begin{synthesis}[the area form of the sphere is a contraction]\label{syn:sphere}
In Section~\ref{sec:tensorfields} the area form of the sphere of radius $r$ was defined by $\alpha_p(v_1,v_2)=\frac1r(v_1\times v_2)\cdot p$. Since $(v_1\times v_2)\cdot N=\vol(N,v_1,v_2)$ for $N=p/r$ (triple product),
$\alpha=\iota_N\vol$ restricted to the sphere: this is Lemma~\ref{lem:fluxbdry} with $X=N$, $\langle N,N\rangle=1$. In spherical coordinates it is $r^2\sin\theta\,\dd\theta\wedge\dd\varphi$, with total area $4\pi r^2$, and Example~\ref{ex:ball} (flux of the position field through $S^2$, $=4\pi$) is the case $r=1$.
\end{synthesis}

% Section 19.5: 05-what-needs-what | Part: 3 | Depends on: sec:stokes, sec:interior, sec:metric | Status: done | Original lines: 2892-2899
%!TEX root = ../../../main.tex
\begin{remark}[What needs what]\label{rem:whatneeds}
\begin{itemize}
\item \emph{Stokes:} smooth structure $+$ orientation only. No metric, no volume.
\item \emph{Divergence/Gauss:} Stokes $+$ a volume form $\mu$ to turn a vector field into an $(n-1)$-form via $\iota_X\mu$. A metric is needed only to rewrite the flux as $\langle X,N\rangle\,\dd A$. On non-orientable manifolds the theorem survives if $\mu$ is a density.
\item \emph{Kelvin--Stokes / curl:} orientation, a metric and dimension $3$, since the identification of $2$-forms with vector fields ($\hodge$) is special to $n=3$. In other dimensions the object is $\dd\omega$, not a ``curl vector''.
\end{itemize}
\end{remark}

% Section 19 (keypoints): Keypoints | Part: 3 | Depends on: sec:stokes, sec:interior, sec:metric | Status: todo | Original lines: 2900-2908
%!TEX root = ../../../main.tex
\begin{keypoints}[Part III]
\begin{itemize}
\item $\int_\gamma\omega=\int\langle\omega,\dot\gamma\rangle\,\dd t$: integration of a $1$-form is contraction with the velocity; $\int_\phi\alpha$ is full contraction with the tangent $k$-frame.
\item Stokes: $\int_M\dd\omega=\int_{\partial M}\omega$, equivalently $\langle c,\dd\alpha\rangle=\langle\partial c,\alpha\rangle$. Metric-free.
\item Gauss is Stokes for $\omega=\iota_X\mu$ with $\dd\iota_X\mu=\Lie_X\mu=\operatorname{div}_\mu X\,\mu$.
\end{itemize}
\end{keypoints}


